\documentclass[11pt]{article}
\usepackage[margin=1in]{geometry}
\usepackage[T1]{fontenc}
\usepackage{lmodern,amsmath,amssymb,amsthm,mathtools}
\usepackage{booktabs,longtable,array,enumitem,microtype}
\usepackage{xcolor,xurl,hyperref}
\hypersetup{colorlinks=true,linkcolor=blue!45!black,urlcolor=blue!55!black,citecolor=blue!45!black}
\setlength{\emergencystretch}{3em}
\setlength{\parskip}{0.35em}
\newtheorem{theorem}{Theorem}[section]
\newtheorem{lemma}[theorem]{Lemma}
\newtheorem{proposition}[theorem]{Proposition}
\newtheorem{corollary}[theorem]{Corollary}
\newcommand{\Z}{\mathbb Z}
\newcommand{\N}{\mathbb N}
\newcommand{\Pos}{\operatorname{Pos}}
\newcommand{\Con}{\operatorname{Con}}
\newcommand{\ZFC}{\mathrm{ZFC}}
\newcommand{\rd}{\mathrm d}
\newcommand{\repo}{\texttt{ProveIt}}
\title{Radical majorants, guard compression, and Pell descent\newline for Diophantine consistency equations}
\author{A research contribution prepared for Vladimir Reshetnikov's ProveIt project\\
\small Mathematical development and executable audits with AI assistance}
\date{23 October 2026}

\begin{document}
\maketitle
\begin{abstract}
We construct four explicit integer polynomials whose solvability is equivalent
to the existence of a finite ZFC refutation, by composing the repository's
ordinary native calibration with new, exactly verified arithmetic reductions.
The resulting unrestricted-integer forms have variable/degree pairs
$(31,6)$, $(26,8)$, $(23,12)$, and $(22,14)$.
In particular the last form replaces the inspected $(26,16)$ frontier by
four fewer variables and a lower degree. A direct polynomial-part majorant
also lowers the degree of the nine-positive-variable construction from
$1{,}379{,}730$ to $1{,}158{,}228$, with a larger circuit. A separate positive-integer
construction preserves eleven variables and reduces degree $325$ to $277$,
using $111$ live binary arithmetic operations.
The central integer idea combines a canonically calibrated determinant
inequality with a universal simultaneous-positivity identity, then chooses
coordinate elimination and polynomial row operations jointly.
The positive construction removes an unnecessary denominator guard by a
rational Pell descent.
All coefficient tuples are fixed and reproduced exactly. The new algebraic
reductions and degree certificates are checked independently. Their
consistency interpretation uses the inherited ordinary native/Pell
calibration and its canonical completeness theorem. This work does not
independently re-establish that complete semantic bridge or supply its
end-to-end Lean/Rocq verification or an internally encoded ZFC derivation.
\end{abstract}
\newpage
\tableofcontents
\newpage

\section{The target and the comparison}
Fix the recursive syntactic proof predicate used by the repository and let
$\Con(\ZFC)$ mean that it has no finite proof of contradiction.
The target for a polynomial $P\in\Z[x_1,\ldots,x_n]$ is
\begin{equation}\label{eq:target}
 \ZFC\vdash
 \bigl(\neg\exists\boldsymbol{x}\in\Z^n\;P(\boldsymbol{x})=0\bigr)
 \longleftrightarrow\Con(\ZFC).
\end{equation}
For the nine- and eleven-variable constructions, replace $\Z^n$ by
$\Z_{>0}^{9}$ and $\Z_{>0}^{11}$, respectively.
Neither this statement nor its proof asserts the truth
of $\Con(\ZFC)$.

The inspected repository state is the full commit
\begin{center}\small\ttfamily
2605669a532342f153b13eb6b4a0af4b54f5fa37.
\end{center}
It includes all wave-ten packages 68--75. Its active review records the
frontier in Table~\ref{tab:frontier}; the table below compares domains and
degrees explicitly. There is no established global minimum for the user's
word ``simplest.'' A lower variable count, lower degree, shorter arithmetic
circuit, and smaller coefficients are distinct objectives.

\begin{table}[htbp]\centering\small
\begin{tabular}{@{}llll@{}}
\toprule
Domain / criterion & Inspected frontier & New construction & Residuals\\
\midrule
Integers, degree 6 & 32 variables & \textbf{31 variables} &18\\
Integers, degree 8 & 27 variables & \textbf{26 variables} &14\\
Integers, degree 12 & no smaller-degree counterpart & \textbf{23 variables} &11\\
Integers, degree at most 16 & 26 variables, degree 16 & \textbf{22 variables, degree 14} &10\\
Positive integers, 11 variables & degree 325, 113 gates & \textbf{degree 277, 111 gates} &2\\
Positive integers, 9 variables & degree 1,379,730 & \textbf{degree 1,158,228} &1\\
\bottomrule
\end{tabular}
\caption{New explicit tradeoffs against the pinned wave-ten review.
Each integer polynomial is the sum of the squares of its listed residuals.
The eleven-positive equation uses a product finalizer; the nine-positive
equation uses a two-root norm. No optimized scalar-circuit
claim is made for the integer forms.}\label{tab:frontier}
\end{table}

The new $(23,12)$ form is smaller in both variables and degree than the old
$(26,16)$ form. It trades degree against the old $(27,8)$ form.
The new $(22,14)$ form improves that older $(26,16)$ point by $4/26$, or
about $15.38\%$, in variable count, while lowering degree from $16$ to $14$. These improvements do not lower the
already available nine-positive or ten-unrestricted-integer variable bounds.

\paragraph{What is new here.}
The determinant inequality is from package 71, the simultaneous-positivity
lemma from package 68, the signed-scale bootstrap from package 75, and the
rational descent from package 69. We do not claim priority for these
ingredients. The new results are their compatible composition, the resulting
literal equations and elimination schedules, and the observation that
reducing an auxiliary norm modulo a retained strong norm avoids an unnecessary
degree increase; centering the remaining root then reduces it further.
The independent polynomial-part majorant is a different simplification
that retains nine positive variables and reduces degree by about $16.05\%$. None of these combinations is in the pinned wave-ten
frontier. No historical novelty outside the inspected sources is asserted.

\section{A complete 23-integer equation}
\label{sec:23}
All $23$ variables in the following ordered tuple range independently over
\emph{all} integers:
\[
 (J,F,\alpha,f,h,r,v,\tau,\eta,y,\gamma,o,j,E,c,a,I,u,Z,t,p_1,p_2,p_3).
\]
The six fixed integer coefficients are
\[
 \beta, K_0, M_C, M_F, A_*=2^{1113},\ B_*=2^{666104}.
\]
Here $\beta=2^{1064921}-1$, and the first four integers are exactly the
whole source-34 tuple specified in Section~\ref{sec:coeff}. In particular
$M_F$ is already shifted. None of these six integers is an unknown.

Define the following \emph{polynomial abbreviations}; they introduce no
quantifiers or extra equations:
\begin{align}
 q&=\beta J+1,& C&=q-F-\alpha,& Y&=(2v+1)q^3,\notag\\
 X&=4qu^2,& H&=4a+3,& \Delta&=a^2+H,\notag\\
 V&=fo-c,& k&=1-V+jc+hE,& \zeta&=k-\eta,\notag\\
 P&=k+E(c-\eta),&&D=X+ac+\gamma H.\label{eq:23-abbrev}
\end{align}
The packing and sign registers are
\begin{gather*}
 R=((q-1)(q-F)+\alpha)(q^2-1)+(M_C+qM_F)J,\\
 U_*=A_*F,\qquad W_*=B_*C,\qquad
 G=U_*W_*-(q-U_*-W_*)^2,\\
 e=\eta\zeta,\qquad
 \Pos(p_1,p_2,p_3)=1+p_1^2+p_2^2+p_3(p_3+1).
\end{gather*}
The eleven residuals are
\begin{align}
 b_1&=a-E-Y, &
 b_2&=EX+E-Xa,\notag\\
 b_3&=c-(a-E)k-\eta,&
 b_4&=\tau^2-P^2+Pk-1,\notag\\
 b_5&=D^2-\Delta c^2-1,&
 b_6&=jc-V-R,\notag\\
 b_7&=(K_0q+X)C-qF-Z(q-1),&&\notag\\
 b_8&=I-(2r+1)c^2,&
 b_9&=I^2-\Delta(f^2-1),\label{eq:23-residuals}\\
 b_{10}&=y^2-I^2(y^2-V^2)-1,&&\notag\\
 b_{11}&=Ge-(t^2-G-e+1)^2-\Pos(p_1,p_2,p_3).&&\notag
\end{align}
The promised single ordinary polynomial is
\begin{equation}\boxed{\quad P_{23}=\sum_{i=1}^{11}b_i^2=0.\quad}\label{eq:P23}\end{equation}
Every exponent is fixed and nonnegative. There are no divisions, variable
exponents, bounded universal quantifiers, or unnamed auxiliary witnesses in
this equation. The variable called \texttt{v\_scale} in the executable packet
is $v$ above.

\begin{theorem}\label{thm:23}
The polynomial $P_{23}$ has exact total degree $12$ after the six fixed
coefficients are specialized to their actual values. Composing the ordinary
source-34 native calibration with the reductions proved below yields
\eqref{eq:target} for $P=P_{23}$.
\end{theorem}
The exact degree assertion does not depend on a generic-coefficient
argument. The coefficient of $\eta^4 f^4o^4$ in $P_{23}$ is exactly $1$.
Indeed the degree-three part of $e$ contains $-\eta fo$; the only residual
that can contribute this monomial after squaring is $b_{11}$, whose relevant
term is $-\eta^2f^2o^2$. Every residual has degree at most six.
The independent exact expansion checks this coefficient and all eleven
compact formulas against the separate elimination construction.

\section{Compressing the sign constraints}
\label{sec:guards}
\subsection{One universal equation for two positive integers}
For unrestricted integers $x,y,t$, put
\[
 \Gamma(x,y,t)=xy-(t^2-x-y+1)^2.
\]
\begin{lemma}[Nearest-square positivity]\label{lem:pair}
For integers $x,y$,
\[
 x>0\ \wedge\ y>0
 \quad\Longleftrightarrow\quad
 \exists t\in\Z\;\Gamma(x,y,t)>0.
\]
\end{lemma}
\begin{proof}
If $\Gamma>0$, then $xy>0$, so $x,y$ have the same nonzero sign. If both
were negative, say $x=-a,y=-b$ with $a,b\ge1$, then
$(t^2+a+b+1)^2>ab=xy$, a contradiction.

Conversely put $N=x+y-1\ge1$. Then $xy-N=(x-1)(y-1)\ge0$.
Choose the square nearest to $N$ as follows. With $m=\lfloor\sqrt N\rfloor$,
take $t=m$ if $N-m^2\le m$, and $t=m+1$ otherwise. If $N$ is a square,
the error is zero. Otherwise $|t^2-N|\le m$ and $m^2<N$.
In either case $(t^2-N)^2<N\le xy$, proving the claim.
\end{proof}
Every value of $\Pos$ is a positive integer. Conversely every positive
integer $n$ is a value of $\Pos$: apply the three-square theorem to
$4n-3$, which is $1$ or $5$ modulo eight. Exactly one square is odd and
two are even, giving
$4n-3=(2p_3+1)^2+(2p_1)^2+(2p_2)^2$.
Therefore
\begin{equation}\label{eq:pairguard}
 x,y>0\quad\Longleftrightarrow\quad
 \exists t,p_1,p_2,p_3\in\Z\;
 xy-(t^2-x-y+1)^2=\Pos(p_1,p_2,p_3).
\end{equation}
This replaces six private variables in two separate $\Pos$ equations by
four variables in one equation. It is an existence equivalence, not an
identity of the old and new witness sets.

\subsection{A determinant with a known canonical orientation}
The source Q41 has the three positive-product conditions
\[
 F\alpha>0,\qquad FC>0,\qquad\eta\zeta>0,
 \qquad C=q-F-\alpha.
\]
Package 71 replaces the first two with $G>0$, for the $G$ in
\eqref{eq:23-abbrev}. The following universal part is short enough to repeat.

\begin{lemma}[Determinant soundness]\label{lem:det}
For arbitrary real $F,C,\alpha,q$ with $q=F+C+\alpha$ and fixed
$A_*,B_*\ge2$, the inequality $G>0$ implies
$F\alpha>0$ and $FC>0$.
\end{lemma}
\begin{proof}
Write $U=A_*F$, $W=B_*C$, and $T=q-U-W$.
The inequality $UW>T^2$ makes $U,W$ nonzero with the same sign, and
$|T|<\sqrt{UW}\le(|U|+|W|)/2$.
Consequently $q=U+W+T$ has their sign and
\[
 |q|>(|U|+|W|)/2\ge |F|+|C|.
\]
It follows that $\alpha=q-F-C$ has this same nonzero sign.
The argument includes negative $q$ and requires no native or Pell assumption.
\end{proof}
Completeness cannot be inferred for arbitrary old positive triples. It uses
the source's canonical witnesses. For the authentic source-34 tuple, package
71 proves uniformly over every canonical accepting word that
\begin{equation}\label{eq:canonicalbox}
 q/2\le U_*<4q/5,\qquad q/2\le W_*<3q/5,\qquad q>0.
\end{equation}
For clarity, put $s=q/2$, $U_*=s+x$, $W_*=s+y$.
Then $0\le x<3s/5$, $0\le y<s/5$ and
\[
 G=s^2+s(x+y)-(x^2+xy+y^2)>\frac{12}{25}s^2>0.
\]
The canonical interval estimates are justified by the fixed leading
native digits; Section~\ref{sec:coeff} records the numerical alignment.
They survive all rank-tuning dummy changes allowed by that source.

\subsection{The new combined guard}
Apply \eqref{eq:pairguard} with $x=G$ and $y=\eta\zeta$.
The single residual $b_{11}=0$ is equivalent to
\[
 G>0\quad\hbox{and}\quad\eta\zeta>0.
\]
Thus every new zero restores all three old sign conditions, without any
prior sign assumption on a source coordinate. Conversely a canonical
source zero satisfies $G>0$ and $\eta\zeta>0$, so the nearest-square
choice of $t$ and a three-square representation supply the four new
coordinates. The original nine private sign coordinates have been replaced
by four. This is the main variable-count saving.

\section{Why the signed-scale and determinant reductions compose}
\label{sec:compatibility}
The Q41 source contains
\[
 X=4q w,\qquad X=2u_{\rm old}^2,\qquad
 (K_0+4w)C=F+z(q-1).
\]
Package 75 proves that one may replace them by
\begin{equation}\label{eq:bootstrap}
 X=4qu^2,\qquad (K_0q+X)C-qF=Z(q-1),
\end{equation}
removing one scale variable. Its proof first treats both possible signs of
$X$ and $q$. It does not infer $X>0$ from $X=4qu^2$.

The inherited signed-bootstrap theorem has these exact two parts:
\begin{enumerate}[leftmargin=*]
\item With the other Q41 arithmetic equations and the three old positive
product conditions, every integer solution of \eqref{eq:bootstrap}
recovers $q>0$, $X=2^R$, and odd $R$ by Pell rank and Mersenne-congruence
arguments. It restores $w=u^2$, $u_{\rm old}=2^{(R-1)/2}$ and $z=Z/q$.
The last quotient is integral by reducing the cleared transport modulo $q$.
\item Canonical refutation witnesses have $q=2^s$ with odd $s$,
odd $R$, and $R>3q$. The assignments
$u=2^{(R-s-2)/2}$ and $Z=qz$ establish \eqref{eq:bootstrap}; all of
$q,F,\alpha,C,\eta,\zeta$ remain unchanged.
\end{enumerate}
The new combined guard supplies the old product conditions before invoking
the soundness theorem. In the reverse direction the signed substitution
does not change the coordinates in the canonical determinant estimates
\eqref{eq:canonicalbox}. Thus there is a single compatible canonical witness
for both constructions. Equiconsistency alone would not establish this
compatibility; the actual coordinate maps do.

For reference, the nontrivial sign argument in package 75 obtains
$|X|\ge4|q|$, $|E|\ge4|q|^4>3R+1$ before fixing a sign. If $X<0$,
Pell growth and the signed rank equality force $|X|=2^R$ and hence
$q=-2^s$. But $q\equiv1\pmod{2^{1064921}-1}$ forbids a negative power of
two: reduce $s$ modulo $1064921$. This excludes that branch. The complete
ordinary proof and its previous independent review are included as inherited
sources, rather than replacing the argument by sample calculations.

\subsection{Exact elimination}
After signed-scale replacement and the combined guard there are $35$
integer variables and $23$ residuals. The following twelve monic definitions
are eliminated successively:
\begin{align*}
 N&=q(L+1),& C_2&=c^2,& F_2&=f^2,& g&=y^2-V^2,\\
 X&=4qu^2,& B_{\rm aux}&=I^2,&
 D&=4qu^2+ac+\gamma(4a+3),&&\\
 \Delta&=a^2+4a+3,&P&=\eta+\zeta+E(c-\eta),&&&&\\
 L&=q^2-1,& \zeta&=1-V+jc+hE-\eta,& V&=fo-c.&&
\end{align*}
Substitution in every remaining row gives exactly
\eqref{eq:23-residuals}. Monic elimination over $\Z$ is reversible:
the eliminated coordinates have the unique displayed integer values.
The archive contains a full parent-row restoration certificate, including
the twelve rows that become identically zero. Therefore this last step is
a genuine bijection of zero sets, unlike the earlier sign-guard changes.

Combining the all-input soundness, canonical completeness, and exact
eliminations proves Theorem~\ref{thm:23} using the named ordinary source
calibration. All newly supplied steps are elementary arithmetic and finite
polynomial reasoning available in ZFC.

\section{A 22-integer equation of degree 14}
\label{sec:22}
An indiscriminate substitution of $I=(2r+1)c^2$ in $P_{23}$ raises the
auxiliary residual to degree ten and the single polynomial to degree twenty.
There is a better order of operations.
Put $W=y^2-V^2$ and replace the auxiliary residual by
\begin{equation}\label{eq:row-op}
 b'_{10}=b_{10}+Wb_9
 =y^2-\Delta(f^2-1)(y^2-V^2)-1.
\end{equation}
Keeping $b_9$, this is an invertible polynomial row operation: subtract
$Wb_9$ to recover $b_{10}$. It therefore preserves the common zero set
over every commutative ring, not only over a finite collection of tests.

Now delete $I$ and $b_8$, substitute $I=(2r+1)c^2$ only in $b_9$, and use
\eqref{eq:row-op}. Finally set $y=V+y_t$, where $y_t$ is a new unrestricted
integer replacing $y$. This polynomial shear has inverse $y_t=y-fo+c$
and changes $y^2-V^2$ to $y_t(2V+y_t)$, of degree three instead of four.
The resulting explicit polynomial is
\begin{equation}\label{eq:P22}
\boxed{\begin{aligned}
 P_{22}={}&\sum_{i=1}^{7}b_i^2+
 \bigl((2r+1)^2c^4-\Delta(f^2-1)\bigr)^2\\
 &+\bigl((V+y_t)^2-\Delta(f^2-1)y_t(2V+y_t)-1\bigr)^2+b_{11}^2=0.
\end{aligned}}
\end{equation}
All abbreviations of Section~\ref{sec:23} are retained except $I$, which
no longer occurs, and $y$ is replaced by $y_t$. There are $22$ unrestricted
integer variables and ten residuals.

\begin{theorem}
The integer zero sets of $P_{22}$ and $P_{23}$ are in bijection under
insertion/removal of $I=(2r+1)c^2$ and the stated root shear. The degree of
$P_{22}$ is exactly $14$.
It satisfies \eqref{eq:target} with the same semantic dependencies as
$P_{23}$.
\end{theorem}
\begin{proof}
The row operation is invertible while $b_9=0$ is retained; the subsequent
monic elimination and polynomial shear are uniquely reversible. A sum of integer squares is zero
exactly when each summand is zero.
The new strong residual has degree six. The centered auxiliary residual has
degree seven, since $V=fo-c$ has degree two. All other residuals have degree
at most six. The coefficient of $a^4f^6o^2y_t^2$ in the final polynomial is
exactly $4$, arising uniquely from the square of the term $-2a^2f^3oy_t$
in the centered auxiliary residual. This proves the exact degree after specialization of all ports.
\end{proof}
The elementary row reduction is useful independently of this application:
one should minimize degree \emph{modulo the retained equations} before
substituting high-degree definitions into every occurrence. Here direct
substitution gives degree $20$, row reduction gives $16$, and the subsequent
root centering gives $14$, at the same $22$ variables.

\section{The sextic and octic companions}
\subsection{Twenty-six unrestricted integers at degree eight}
Use the ordered variables
\[
 (J,F,\alpha,z,f,h,r,v,w,\tau,\eta,\zeta,y,\gamma,o,j,E,c,a,V,I,u,t,p_1,p_2,p_3).
\]
Keep $q,C,H,\Delta,R,G,e,\Pos$ as before, but now set
\[
 X=2u^2,\quad Y=(2v+1)q^3,\quad k=\eta+\zeta,\quad
 P=k+E(c-\eta),\quad D=X+ac+\gamma H.
\]
The fourteen residuals are
\begin{align*}
 &X-4qw,\quad a-E-Y,\quad EX+E-Xa,\quad c-(a-E)k-\eta,\\
 &\tau^2-P^2+Pk-1,\quad D^2-\Delta c^2-1,\quad k-hE-R-1,\\
 &(K_0+4w)C-F-z(q-1),\quad V+c-fo,\quad V-cj-hE+k-1,\\
 &I-(2r+1)c^2,\quad I^2-\Delta(f^2-1),\quad
 y^2-I^2(y^2-V^2)-1,\\
 &Ge-(t^2-G-e+1)^2-\Pos(p_1,p_2,p_3).
\end{align*}
Their sum of squares is $P_{26}$. Every residual has degree at most four.
Only the last uses $t$, with coefficient $-1$ at $t^4$, so
$[t^8]P_{26}=1$. Keeping $\zeta$ independent is deliberate: it keeps
$e=\eta\zeta$ quadratic and the joint guard quartic.

Soundness and canonical completeness use Section~\ref{sec:guards}, followed
by ten exact monic eliminations from the unswitched Q41 source. No
signed-bootstrap theorem is needed for this companion. It has $26$
unrestricted integer variables, exact degree eight, and the same consistency
interpretation.

\subsection{Thirty-one unrestricted integers at degree six}
Start with the signed-bootstrap source and replace only the $F\alpha,FC$
guards by $G=\Pos(p_1,p_2,p_3)$. Retain the separate
$\eta\zeta=\Pos(e_1,e_2,e_3)$ row. The variables are
\[
\begin{split}
 &(J,F,\alpha,f,h,r,v,\sigma,\eta,\zeta,y,\gamma,o,j,E,c,a,D,L,V,I,\mathcal B,\Delta,u,\\
 &\hspace{22mm}e_1,e_2,e_3,Z,p_1,p_2,p_3).
\end{split}
\]
Define $q=\beta J+1$, $C=q-F-\alpha$, $k=\eta+\zeta$,
$H=4a+3$, $X=D-ac-\gamma H$, $Y=(2v+1)q(L+1)$, and
\[
 R_L=((q-1)(q-F)+\alpha)L+(M_C+qM_F)J.
\]
The eighteen residuals are
\begin{align*}
 &L-q^2+1,\quad a-E-Y,\quad EX+E-Xa,\quad c-(a-E)k-\eta,\\
 &\sigma^2+(2\sigma-k)E(c-\eta)-1,\quad\Delta-a^2-4a-3,\\
 &D^2-\Delta c^2-1,\quad k-hE-R_L-1,\quad
 (K_0q+X)C-qF-Z(q-1),\\
 &V+c-fo,\quad V-cj-hE+k-1,\quad I-(2r+1)c^2,\\
 &\mathcal B-I^2,\quad\mathcal B-\Delta(f^2-1),\quad
 y^2-\mathcal B(y^2-V^2)-1,\\
 &X-4qu^2,\quad\eta\zeta-\Pos(e_1,e_2,e_3),\quad
 G-\Pos(p_1,p_2,p_3).
\end{align*}
Here $G$ is still the fixed determinant expression in $q,F,C$.
Their sum of squares is $P_{31}$, of exact degree six: the residual
$I-(2r+1)c^2$ contributes the unique coefficient $4$ at $r^2c^4$.
The centered coordinate is $\tau_{\rm old}=\sigma+E(c-\eta)$; this
triangular substitution removes the potentially quartic term from the
first norm before eliminating its product register. Six monic eliminations
then leave the eighteen displayed rows.

\section{Fixed coefficients and exact provenance}
\label{sec:coeff}
The integer family uses one coherent native tuple, not a mixture of small
ports from different compilers. Its basic quantities are
\[
 \rho=2^{121},\qquad B=\rho^{8801}=2^{1064921},\qquad\beta=B-1.
\]
The shipped literal machine table has $430$ transitions and starts in state
$0$. The deterministic coefficient constructor rebuilds the ordered $943$
tile list from that table, compares every tile with the retained literal
list, and computes the native ports by finite sums and shifts. It performs
no unbounded search for a machine halt or a refutation.

For a compact exact description of those finite sums, let the ordered tiles
be $T_i=(\ell_{i0},\ell_{i1},\ell_{i2},\ell_{i3})$, $0\le i<943$.
Set $w_j=2^{10+2j}$. Define a coefficient map $c_e$ by
\begin{align*}
 c_i&=1+\sum_{s=0}^3\sum_{j=0}^9
        \operatorname{bit}_j(\ell_{is}+512)w_{10s+j}
        +\sum_{j=40}^{43}w_j,\qquad 0\le i<943,\\
 c_{o_s+2j}&=w_{10s+j},\quad(o_0,o_1,o_2,o_3)=(1020,1000,980,960),\\
 c_{1095}&=w_{40},\quad c_{1096}=w_{41},\quad
 c_{2195}=w_{42},\quad c_{3296}=w_{43},\quad c_{1063}=c_{1064}=0.
\end{align*}
Let $\mathcal E$ be these $989$ support positions, including the two zero
controller positions. Then
\begin{align*}
 D_C&=\rho^{40}+\rho^{41}+\rho^{1100}+\rho^{2200}
       +\sum_{e\in\mathcal E}c_e(\rho^{4393-e}+\rho^{5495-e}),\\
 K_0&=D_C+B\rho,\qquad
 M_C=B-1-\sum_{e\in\mathcal E\setminus\{1\}}\rho^e-2\rho^{1063},\\
 \mu&=1022+\sum_{j=0}^{43}w_j,\\
 M_F&=B-1+4+\mu(\rho^{4393}+\rho^{5495})
       +\sum_{j=0}^9(\rho^{1000+2j}+\rho^{1021+2j})
       +\rho^{2195}+\rho^{3296}\\
    &\hspace{8mm}+(2^{119}-1)\sum_{e=8792}^{8798}\rho^e
       +(2^{27}-1)\rho^{8799}.
\end{align*}
All exponents in this fixed recipe are nonnegative. The finite tile list,
the machine table generating it, and all bytes of the six resulting integers
are included. They are alternative exact specifications of fixed integers;
the hashes in the manifest are checks, not definitions of the coefficients.

The determinant alignment uses the leading coefficient
$c_0=105534389824824194647382298625$ and the identities
\[
 121\cdot8791+1113=1064921-97,\qquad
 121\cdot3296+666104=1064921-1.
\]
Canonical native blocks have $C$'s leading digit $1$ at $3296$ and $F$'s
leading coefficient $c_0$ at $8791$. Together with
$2^{96}\le c_0<c_0+1\le3\cdot2^{95}$ and the source's carry bounds,
these yield \eqref{eq:canonicalbox}. The retained determinant proof gives
the full finite-to-canonical argument.

\begin{center}
\begin{tabular}{@{}lrr@{}}
\toprule Fixed integer & Bit length & Minimal unsigned bytes\\\midrule
$\beta$ &1064921&133116\\
$K_0$ &1065043&133131\\
$M_C$ &1064921&133116\\
$M_F$ &1064922&133116\\
$A_*$ &1114&140\\
$B_*$ &666105&83264\\\midrule
Total &&615883\\\bottomrule
\end{tabular}
\end{center}
These coefficient costs are unchanged from package 71. The new savings are
in variables and degree. The eleven-positive construction uses the distinct
whole compact44 tuple with $\beta=2^{590069}-1$; it has its own binding and
does not share coefficient leaves with the integer family.
The nine-positive construction uses the distinct whole original31 tuple.
Although its $\beta$ is also $2^{1064921}-1$, its other three ports differ
from source34. Its raw four-port payload is $533099$ bytes.
Equality of the radix does not authorize mixing these tuples.

\section{Eleven positive variables: degree 277}
\label{sec:positive}
Package 70 combines three square roots into one norm and uses the strong
Pell ordinate $I=i q(q-1)c^2$. The factors $q(q-1)$ ensure a coprimality
condition used only to recover integrality of an auxiliary rational root.
Package 69's following stronger lemma makes that condition unnecessary.

\begin{lemma}[Rational Pell descent]\label{lem:descent}
Let $Q>1$ and $V$ be integers, and $y\in\mathbb Q$. If
$QV^2-(Q-1)y^2=1$, then $y\in\Z$.
\end{lemma}
\begin{proof}
Neither coordinate is zero; change signs to make them positive. Write
$y=a/m$ in lowest terms, $m>0$. The equation gives $m^2\mid Q-1$.
If $V=1$, then $y=1$. For $V>1$ define
\[
 V'=(2Q-1)V-2(Q-1)y,\qquad y'=(2Q-1)y-2QV.
\]
The first is integral, the norm is unchanged, and $y'$ has the same reduced
denominator $m$ because $2Q-1\equiv1\pmod m$.
The inequalities $y>V$ and $QV>(Q-1)y$ follow from the original norm.
They imply $-V<V'<V$, while the new norm excludes $V'=0$.
Replace the new coordinates by absolute values and induct on the positive
integer $V$. The denominator is preserved until $V=1$, proving $m=1$.
\end{proof}

The new equation simply replaces $I$ by $ic^2$ throughout package 70's
literal norm. All unconditional radical separation and transport domination
estimates need only $I\ge c^2$, which remains true for positive $i$.
Once the norm recovers rational square roots, two are integral as rational
roots of integer squares, and Lemma~\ref{lem:descent} makes the third
integral. This is the only step at which the removed factor was used.

For an exact standalone polynomial specification, the positive unknowns are
\[
 (j,F,\alpha,i,T,X,Y,n,\eta,\zeta,\gamma).
\]
With the compact44 coefficients, define
\begin{align*}
 &J=3j+1,\quad q=\beta J+1,\quad C=q-F-\alpha,\quad E=XY,\\
 &k=\eta+\zeta,\quad c=kY+\eta,\quad a=E+Y,\quad H=4a+3,\quad
 \Delta=a^2+H,\quad\ell=EkY,\\
 &R=((q-1)(q-F)+\alpha)(q^2-1)+(M_C+qM_F)J,\quad
 \widehat R=q^3(q+F+\alpha),\\
 &M=q^3(q-1),\quad B_t=(qK_0+X)C-qF,\quad D_\tau=\ell(\ell+k)+1,\\
 &I=ic^2,\quad D_f=1+\Delta I^2,\quad Q=\Delta^2I^2,\quad\delta=Q-1,\\
 &V=D_f(cT-R)-c,\quad D_n=QV^2-1,\quad\overline V=D_f(cT+\widehat R),\\
 &L=3\overline V,\quad W=1+(q-1)L,\quad
 \overline Z=M(\ell+k+n+L)+(q-1)Y-q^3WB_t,\\
 &r_0=k-R-1,\quad S=4qY,\quad U=qS,\\
 &D_*=U(X+ac)+H(U\gamma+SX-r_0),\quad
 A_1=D_*^2-U^2(1+\Delta c^2),\\
 &A=M^2D_\tau,\quad B_f=M^2D_f,\quad C_n=M^2D_n,\\
 &T_n=\delta^2AB_f,\quad R_n=\delta C_n,\quad
 J_n=\delta(\overline Z^2-A-B_f)-C_n,\\
 &N_c=J_n^2-4\bigl(T_n+R_n(A+B_f)\bigr),\quad
 N_8=N_c^2-64T_nR_n\overline Z^2.
\end{align*}
Then
\begin{equation}\boxed{\quad P_{11}=N_8+A_1+N_8A_1=0.\quad}\label{eq:positive11}\end{equation}
The full proof of both directions, including signed radical choices,
positive reconstruction of the deleted transport gap, the product finalizer,
and exact degree, is included in \texttt{positive11/PROOF.md}.

Briefly, the eight-root norm gives a rational signed sum of square roots of three positive,
distinct radicands, with no zero signed sum. Hence all three roots are
rational. The two integer radicands and Lemma~\ref{lem:descent} make every
root integral. The unchanged bound $L>|f_*|+|y_*|$ restores the positive
transport gap $r=n+L-f_*-y_*$, and $|y_*|>|V|$ restores the positive
auxiliary gap. In the completeness direction canonical Pell witnesses
already permit $c^2\mid I$, so the new parameter $i=I/c^2$ is positive.
One chooses the weighted transport gap after the strong and auxiliary
roots; its weight $L$ is large enough for the final $n$ to be positive.

The product finalizer is exact: $U=4q^2Y$ gives
$A_1\equiv D_*^2\pmod4$, and the factor $64$ gives
$N_8\equiv N_c^2\pmod4$. Thus neither residual is $-2$; no parity
assumption on $q$ is required. Thus
$(1+A_1)(1+N_8)=1$ forces both factors to be $+1$, and both residuals to
vanish. Its exact degree is $277$, certified by the nonzero coefficient
\[
 [j^{60}i^{24}\eta^{49}X^{42}Y^{85}\alpha^{16}\gamma]P_{11}
 =2^7\,3^8(3\beta)^{60}>0.
\]
The literal division-free circuit uses $111$ live binary additions,
subtractions and multiplications, including scalar multiplications and
squaring. Fixed integer leaves and register reuse are free in this metric.
The two-residual system itself uses $108$ operations. The degree reduction
is $48/325\approx14.77\%$.

\section{A direct majorant of the complete Pell radical}
\label{sec:directmajorant}

Another new construction improves a different measure of simplicity:
the total degree with nine positive unknowns.  It retains package~71's
source encoding and replaces its final square-root bound.  The resulting
polynomial has degree $1\,158\,228$, compared with $1\,379\,730$ for that
package, and its explicit circuit has $1\,753$ arithmetic operations,
compared with $478$.  Thus this is a degree--circuit-size tradeoff.
The improvement is approximately $16.05\%$ in degree.  It neither lowers
the number of unknowns nor asserts a global minimum.

\subsection{Why approximate the complete radical?}

The previous finalizer constructs a polynomial upper bound for
$\sqrt{DF}\,G$ and then multiplies it by $H$.  Errors in that bound are
therefore multiplied by a high-degree polynomial.  We instead construct
a polynomial upper bound for the complete radical
\[
 \sqrt{DF\bigl((G^2-1)H^2+1\bigr)}.
\]
The discovery calculation uses a formal expansion at $A=\infty$.
The proof of the resulting bound consists entirely of finite integer
polynomial identities and positive coefficients.

In this subsection $A,C,i,Z$ are independent variables.  Define
\begin{align*}
 J&=4A-5,& D&=(A^2-1)C^2+1,\\
 e&=iJC,& E&=eDC,\\
 F&=(A^2-1)E^2+1,& L_0&=CD+1-A,\\
 G&=A+FL_0,& H&=C+FZ,\\
 I&=(G^2-1)H^2+1,& Q_*&=DFI.
\end{align*}
As a polynomial in $A$, $Q_*$ has degree $46$ and leading coefficient
\[
 (1024C^{24}i^5Z)^2.
\]
Write $Q_*=\sum_{k=0}^{46}q_k A^k$.  Define $t_{23},\ldots,t_0$ in
$\mathbb Q(C,i,Z)$ by the finite recurrence
\begin{align}
 t_{23}&=1024C^{24}i^5Z,\nonumber\\
 t_{23-k}
 &=\frac{q_{46-k}-\sum_{\ell=1}^{k-1}
                 t_{23-\ell}t_{23-k+\ell}}{2t_{23}}
       \qquad(1\leq k\leq23).                 \label{eq:majorantrecurrence}
\end{align}
This is the polynomial part of the formal square root with positive
leading term.  The recurrence is a finite specification of coefficients;
no division by an unknown remains in the final equation.
Set
\begin{equation}
 \mathcal P(A,C,i,Z)=65536C^7\sum_{r=0}^{23}t_r A^r,
 \qquad \sigma=65536C^7.                       \label{eq:majorantP}
\end{equation}
Exact coefficient arithmetic proves
$\mathcal P\in\mathbb Z[A,C,i,Z]$.  It has $654$ nonzero terms, all of
which are supplied in \texttt{nine\_positive/data/radical.json}.
The same file gives every coefficient of
\begin{equation}
 \mathcal D=\mathcal P^2-\sigma^2Q_*.
                                                    \label{eq:localdefect}
\end{equation}
The coordinate degrees of these two polynomials are
\[
 \deg_{(A,C,i,Z)}\mathcal P=(23,31,5,1),\qquad
 \deg_{(A,C,i,Z)}\mathcal D=(22,42,10,2).
\]
There are $3781$ nonzero defect terms.  Denominators in the formal
coefficient recurrence cancel after the displayed factor $65536C^7$;
there are no negative exponents in either delivered table.

\subsection{A uniform inequality certified by positive coefficients}

The decisive certificates are the following exact identities, represented
by their complete coefficient lists:
\begin{enumerate}
 \item $\mathcal P(X+3,Y+1,i,Z)$ has $1816$ nonzero coefficients,
       all strictly positive.
 \item $\mathcal D(X+3,Y+1,U+1,V+1)$ has $32\,637$ nonzero
       coefficients, all strictly positive.  Its constant coefficient is
       \[
          295870595720775501924817.
       \]
\end{enumerate}
The lists are in \texttt{certificates/P\_shifted.json} and
\texttt{certificates/defect\_shifted.json} within the nine-positive
component.  Regeneration by rational coefficient arithmetic and an
independent implementation using only integer sparse polynomials agree.
These are finite algebraic certificates for statements over the whole
domain, rather than sampled inequalities.

\begin{lemma}[Complete-radical majorant]
For all integers $A\geq3$ and $C,i,Z\geq1$,
\begin{equation}
 \mathcal P(A,C,i,Z)>65536C^7\sqrt{DFI}>0.
                                                    \label{eq:majorantineq}
\end{equation}
\end{lemma}
\begin{proof}
Here $J,D,E,F$ are positive.  Since
$CD+1-A\geq A^2-A+1>0$, we have $G>A\geq3$; also $H,I>0$.
The first certificate evaluated at $X=A-3$, $Y=C-1$ gives
$\mathcal P>0$.  All its monomials are nonnegative, and its terms of
$X$- and $Y$-degree zero are positive when $i,Z\geq1$.
The second certificate evaluated at
$(X,Y,U,V)=(A-3,C-1,i-1,Z-1)$ gives $\mathcal D>0$.
Strictness holds even on every boundary because its constant is positive.
Equation~\eqref{eq:localdefect}, together with $\mathcal P>0$, now gives
\eqref{eq:majorantineq}.
\end{proof}

The use of a formal expansion in the discovery process therefore does
not introduce an asymptotic qualification into the theorem.  The bound
holds before any solution of the final equation has been assumed.

\subsection{The fixed equation and its nine unknowns}

All nine unknowns
\[
 b,g,h,s,w,C,i,j,n\in\mathbb Z_{>0}
\]
retain the domain and meaning in the copied source-71 prefix.  The source
polynomials $R,N,t$ are the fixed source-65 encoding polynomials; their
complete weights, target positions and coefficient kernels are included
in the component's data directory.  The whole original-31 native
coefficient tuple is supplied by literal payloads and a finite generator.
No coefficient is left as a free program parameter.

For clarity, the ratio registers needed below are
\begin{align*}
 U&=N^2w,&Y&=N^2s,& M&=RY,\\
 P_0&=2M^2U,&K&=R+1+h(P_0-1),&A&=M(U+1),\\
 B_1&=2R+1,&W_1&=bw,&J&=4A-5,\\
 Q_0&=3W_1C-2(W_1^2-1),&&&&\\[-6pt]
 m&=K^2(4t-g)-16t(C-KY)^2,&&&&\\
 \Omega&=Jm,&E_0&=J(n+W_1^2)+Q_0,&V_0&=P_0K,\\
 a_0&=(P_0^2-1)K^2+1,& Z&=B_1+(j-1)C.&&
\end{align*}
In the local polynomials above, substitute this $A$, this $C$, the free
$i$, and this $Z$.  In particular $H=C+FZ$ is exactly the inherited
Lucas-certificate register.  Define
\begin{equation}
 V_1=\mathcal P(A,C,i,Z),\qquad
 a_1=(65536C^7)^2DFI,\qquad
 \Theta=E_0-\Omega(V_0+V_1).                 \label{eq:newnineinputs}
\end{equation}
The new equation is
\begin{equation}
 \boxed{\quad
 F_9=(\Theta^2-\Omega^2(a_0+a_1))^2
                       -4\Omega^4a_0a_1=0.\quad}   \label{eq:newnine}
\end{equation}
Every named register is a polynomial abbreviation, not an additional
unknown.  The division-free DAG \texttt{equation9.json} is the complete
literal instance.  The source bounds $R,N\geq8$ imply $A\geq3$ on every
ambient positive input, and $C,i,Z\geq1$ is immediate.  Consequently the
majorant lemma applies without a circular appeal to the final equation.

\subsection{Why the replacement preserves the ratio theorem}

The original first majorant satisfies $V_0>\sqrt{a_0}>0$, since
$V_0^2-a_0=K^2-1>0$.  The lemma gives
$V_1>\sqrt{a_1}>0$.  Moreover
\[
 E_0=Jn+(J-2)W_1^2+3W_1C+2>0.
\]
Thus a zero of~\eqref{eq:newnine} cannot have $m=0$: that would give
$F_9=E_0^4>0$.  Factoring the two-root norm over the reals gives signed
roots $r_0^2=a_0$, $r_1^2=a_1$ with
\begin{equation}
 \frac{E_0}{J}=m(V_0+V_1+r_0+r_1).             \label{eq:marginorientation}
\end{equation}
The strict majorants make the parenthesis positive for every sign choice,
so $m>0$ follows before any root-integrality argument.  The unchanged
margin argument now recovers the required free-carrier bounds and makes
$DFI>a_0$, hence $a_1>a_0$.

The sum $r_0+r_1=\Theta/\Omega$ is rational and nonzero.  The identity
\[
 2(r_0+r_1)r_0=(r_0+r_1)^2+a_0-a_1
\]
therefore makes both roots rational, and an integer with a rational square
root has an integral square root.  Although the scale
$65536C^7$ depends on the current positive input, it is a nonzero integer.
Consequently $r_1/(65536C^7)$ is a rational square root of $DFI$ and is
also integral.  No fixed-scale assumption is needed.  Equation
\eqref{eq:marginorientation} then recovers $J\mid Q_0$ exactly as before.
All the inherited ratio conditions, and hence the source zero, follow.

Conversely, a source zero supplies the original ratio witnesses and
integral roots.  Choose the positive roots
$r_0=\sqrt{a_0}$ and $r_1=65536C^7\sqrt{DFI}$, and take
\[
 n=m(V_0+V_1+r_0+r_1)-W_1^2-Q_0/J.
\]
This is an integer.  The unchanged estimate $W_1^2+Q_0/J<2D$, together
with $F>D$, $I\geq1$ and $V_1>65536\sqrt{DFI}>2D$, shows $n>0$.
One norm factor vanishes.  This proves the same source-zero equivalence
as source~71.  Composing with the retained native-source and syntactic
refutation calibration, and negating, gives the ordinary theorem
formalizable in ZFC:
\[
 \mathsf{ZFC}\vdash
 \bigl(\text{$F_9$ has no positive integer zero}\bigr)
       \ \longleftrightarrow\ \operatorname{Con}(\mathsf{ZFC}).
\]
The new component supplies no additional end-to-end Lean or
object-language ZFC certificate; the inherited calibration boundary is
stated explicitly in its proof and provenance files.

\subsection{Exact degree and the change of dominant term}

Let $r=7963$, $\nu=2115$, $W=116$ be the degrees of $R$, $N$ and the
maximum source weight, and put $a=r+4\nu+2=16425$.  Then
\begin{align*}
 \deg\Omega&=5a+2W+3=82360,&\deg V_0&=3a+r+1=57239,\\
 \deg\delta_0&=3a+r+2=57240,&
 \delta_0&=V_0^2-a_0=K^2-1.
\end{align*}
The unique highest weighted-degree term of $\mathcal P$ is
$2^{26}A^{23}C^{31}i^5Z$, giving
\[
 \deg V_1=23a+r+36=385774.
\]
The corresponding unique term of its defect is
\[
 141664530202624\,A^{22}C^{42}i^{10}Z^2,
\]
so $\delta_1=V_1^2-a_1$ has degree $22a+2r+52=377328$.
To resolve the cancellations in the outer norm, set
\[
 B_*=E_0^2-2E_0\Omega(V_0+V_1)+\Omega^2(\delta_0+\delta_1).
\]
It has degree $542048$, and the exact identity
\begin{equation}
 F_9=B_*^2+4\Omega^2V_0V_1B_*
                 +4\Omega^4(\delta_0V_1^2+a_0\delta_1)       \label{eq:resolvednorm}
\end{equation}
shows the following four degrees:
\begin{center}
\begin{tabular}{lr}
$B_*^2$ & $1\,084\,096$\\
$4\Omega^2V_0V_1B_*$ & $1\,149\,781$\\
$4\Omega^4\delta_0V_1^2$ & $1\,158\,228$\\
$4\Omega^4a_0\delta_1$ & $821\,246$
\end{tabular}
\end{center}
The third term is uniquely largest.  Its leading coefficient is positive,
so the exact degree is
\[
 \boxed{\deg F_9=72r+276\nu+8W+224=1\,158\,228.}
\]
The dominant term is different from that of the previous finalizer;
merely substituting the new defect degree into its old formula would be
incorrect.  If $\rho,\sigma_N,\tau$ are the positive diagonal leading
coefficients of $R,N,t$, respectively, the new leading coefficient is
$2^{90}\tau^4\rho^{72}\sigma_N^{276}>0$.  For the actual source constants
$\alpha=2^{2134939}$ and $c=2^{17079501}$ this becomes
\[
 2^{162}\alpha^{578828}c^{276}>0.
\]
Thus specialization to the actual native coefficients cannot remove the
leading term.

The final circuit contains $923$ multiplications, $810$ additions and
$20$ subtractions.  This count includes every multiplication by a fixed
integer and counts reuse once.  The independent verifier reconstructs
all $654$ majorant coefficients directly from the final DAG, checks the
whole finalizer as a polynomial identity, and verifies $45$ retained
named source expressions against source~71.  Both independent checkers
use only the Python standard library and pass with and without Python's
optimization flag.


\section{Proof boundary and reproducibility}
\label{sec:boundary}
The new transformations are exact ordinary mathematical results. The
semantic dependency is the ordinary native/Pell calibration presented in
the cited repository manuscripts, including all-input soundness and
canonical completeness for the exact coefficient tuple. The present work
proves new reductions using that theorem; it does not independently close
the repository's remaining formal native-calibration or internal-ZFC
obligations. The unbounded semantic input states that
the specified source-34, compact44, or original31 system has a solution if and only
if the fixed calculus has a finite ZFC refutation, with the documented
canonical completeness witnesses. The proofs above use its all-input
soundness and its stronger canonical completeness, with the literal
coefficient bindings. They do not infer either from finite simulations.

Every new integer zero first determines the old product guards and signed
arithmetic source, which in turn determines a finite ZFC refutation.
In the other direction one regenerates a canonical source witness,
preserves its determinant estimates, supplies a nearest-square guard, and
applies the exact coordinate maps. These ordinary arithmetic arguments
can be formalized in ZFC. Both the full formal native/compiler calibration
and a machine-checked proof of the displayed object-language sentence are
separate deliverables, and are not supplied here.

In particular the repository's \texttt{OBJECT\_ZFC.md} correctly distinguishes
an external Lean equivalence from a proof in the sealed ZFC calculus.
Truth on external natural numbers, numeralwise agreement, and bounded
consistency proofs alone do not settle nonstandard internal inputs.
Nothing in this report treats those weaker statements as the requested
internal biconditional. The local polynomial row reductions, monic maps,
and sign lemmas can be inserted directly into that formalization once its
formal native or direct-checker endpoint is complete. This is a distinction
between ordinary theorem dependencies and their formal verification, not an
inference from finite replay to an unbounded consistency statement.

\paragraph{Executable evidence.}
The archive includes literal expanded residual packets, compact formulas,
exact coefficient payloads and regeneration code, source pins and licences,
independent algebraic checks, a top-level verifier, and this LaTeX source.
The integer checker expands over symbolic fixed ports, verifies every
elimination and parent restoration, and exhibits coefficients that remain
nonzero on the actual fixed slice. The independently written audit checks
the final formulas without importing the construction code. Normal and
optimized Python executions retain all explicit checks; finite controls
are labelled as controls, not as proofs of unbounded claims.

\paragraph{Useful next formalization tasks.}
The nearest-square lemma, determinant implication, canonical-interval
composition, and elementary norm-row operation are small reusable targets.
The last has an especially simple ring certificate. These can improve the
closed candidate while the larger native-calibration and internal-checker
projects proceed. The new low-degree forms do not depend on Jones or Sun
compression; their inherited source still uses the K215 native encoding.

\appendix
\section{An alternative route: one loop for unordered proof certificates}
\label{sec:single-loop}

This section gives a source-level improvement independent of Turing-machine
simulation and Jones compression.  It does not claim a new measured final
polynomial.  The affine-product construction in package~68 uses one bounded
universal block for proof steps and a second for a common-multiple condition.
Padding the proof makes the two blocks share one index.  The resulting source
has four outer natural parameters instead of five and one bounded universal
block instead of two.  The three separate label bounds are merged into one joint bound, reducing
its local witness count from $23$ to $22$.
The number of witnesses after arithmetic closure has not been evaluated.

\subsection{Statement}

Let a Horn calculus have natural-number labels, a freely derivable truth
label~$0$, and rules $(p,q)\longrightarrow c$ with one conclusion and two
premise slots.  Axioms have premises $(0,0)$; unary rules are padded by~$0$.
Include the harmless rule $(0,0)\longrightarrow0$ and fix target~$1$.
Write $D$ for its least set of derivable labels.

For $M\geq1$, put $f_M(j)=1+jM$.  Let
\[
 \beta(x,y,i)=x\bmod\bigl(1+(i+1)y\bigr).
\]
The following predicate has outer natural parameters $T,x,y,M_0$, where
$M=M_0+1$:
\begin{align}
 \beta(x,y,0)&=M+1,& \beta(x,y,T)&=1, \label{eq:sl-endpoint}\\
 \forall i<T\quad\exists p,q,c,d,k\in\mathbb N:\quad
 (p,q)&\longrightarrow c, \nonumber\\
 p+q+c+d&=T,\label{eq:sl-bounds}\\
 M&=(i+1)k,\label{eq:sl-common}\\
 \beta(x,y,i+1)f_M(c)&=
 \beta(x,y,i)f_M(p)f_M(q).\label{eq:sl-transition}
\end{align}
The rule relation in this display can have its own fixed local witnesses.
Only one occurrence of bounded universal quantification is present.

\paragraph{Theorem.}
The predicate has a solution if and only if $1\in D$.
This equivalence is provable by elementary finite mathematics in ZFC.

\subsection{Soundness}

The endpoints imply $T>0$: at $T=0$ the same beta value would equal both
$M+1\geq2$ and~$1$.  Equation~\eqref{eq:sl-common}, as $i$ runs from~$0$
to~$T-1$, says that every positive integer at most~$T$ divides~$M$.
For distinct $1\leq a<b\leq T$, the weights $f_M(a)$ and $f_M(b)$ are
coprime.  Indeed, a common divisor $d$ divides $(b-a)M$ and is coprime
to~$M$, hence divides $b-a$.  But $b-a\mid M$, forcing $d=1$.

Put $S_i=\beta(x,y,i)$.  Since all the weights are positive, $S_0=M+1>0$
and~\eqref{eq:sl-transition} imply $S_i>0$ for every $i\leq T$.
Multiplying the transition identities and cancelling the positive intermediate
states gives
\[
 \prod_{i<T} f_M(c_i)
 =f_M(1)\prod_{i<T}f_M(p_i)f_M(q_i).
\]
Pairwise coprimality of the positive-label weights, together with $f_M(0)=1$,
forces exact equality of positive-label multiplicities.  For example, choose
one prime divisor of each $f_M(j)$ and compare its valuation; that divisor
occurs in no other weight.  In the free abelian group on positive labels,
\[
 \sum_{i<T}(e_{c_i}-e_{p_i}-e_{q_i})=e_1,
 \qquad e_0:=0.
\]
If $1\notin D$, sum the finitely many coordinates outside~$D$.  A rule with
conclusion outside~$D$ has at least one premise outside~$D$, while a rule
with conclusion inside~$D$ contributes at most zero.  Thus the left side
has nonpositive sum and the right side has sum~$1$, a contradiction.

This argument does not accept cyclic proofs.  Its use of one conclusion per
rule is essential; the corresponding assertion for several conclusions is
false.

\subsection{Completeness and the order of the choices}

Take a finite proof tree of label~$1$.  Expand it from its root, recording a
rule every time an open nonzero obligation is replaced by its two premises;
zero premises are discarded.  Let $L$ be the number of recorded rules and
$m$ the greatest sum $p+q+c$ in a recorded rule.  Choose
\[
 T\geq\max(1,L,m),\qquad M=T!,\qquad M_0=M-1.
\]
Pad the list to length~$T$ with $(0,0)\longrightarrow0$.  These padding steps
leave the accumulator unchanged because $f_M(0)=1$.

At each stage let $S_i$ be the product of $f_M(j)$ over the current multiset
of nonzero open obligations.  Removing the chosen conclusion and inserting
its premises proves~\eqref{eq:sl-transition}.  The initial value is $M+1$;
the final value is~$1$.  Every row-label sum is at most~$T$, supplying the joint slack; and the quotients
$k_i=M/(i+1)$ establish~\eqref{eq:sl-common}.

Only after these finite integers have been chosen do we choose a beta scale.
For instance, take
\[
 y=T!\bigl(1+\max_{0\leq i\leq T}S_i\bigr).
\]
Then the moduli $1+(i+1)y$, for $0\leq i\leq T$, are pairwise coprime:
their index differences are at most~$T$ and divide~$y$.  They also exceed
the required residues.  The finite Chinese remainder theorem supplies
$x$ with $\beta(x,y,i)=S_i$ simultaneously.  This completes the proof.
The scale~$y$ and the common multiple~$M$ are distinct witnesses; identifying
them would generally destroy completeness.  The arithmetic closure
compiler's own scales must likewise remain independent.

\subsection{Literal local-polynomial adapter}

Package~68 supplies a nonnegative local polynomial
$H(S,S',m,M_0;\mathbf u)$ with $23$ natural witnesses and degree~$50$.
It checks one padded numeric C222 rule, its three label bounds, and its
affine accumulator update.  Write $H_0$ for its Horn product plus accumulator-update square, removing
the three old bound squares and their slack coordinates.  Here $\mathbf u$
denotes the twenty retained witnesses and includes the three label coordinates
$p,q,v$.  With a fresh joint-bound witness $b_0$, define
\[
 J(S,S',T,M_0,i;\mathbf u,b_0,k)
 =H_0(S,S',M_0;\mathbf u)
  + (p+q+v+b_0-T)^2 +\bigl(M_0+1-(i+1)k\bigr)^2.
\]
The retained nonnegativity makes $J=0$ precisely the conjunction of the rule/update, joint-bound, and divisor
conditions.  This source has five exposed step coordinates and $22$ local
natural witnesses.  Its degree remains exactly~$50$: the new summands
have degrees at most~$4$ and cannot remove the degree-$50$ part of the Horn product.

One now applies the supplied bounded-forall constructor once to the single
family $J=0$, conjoins the beta endpoint conditions, and existentially binds
the four outer parameters.  This is a precise source recipe, but the full
ordinary polynomial and its total cost have not been materialized here.
The old C222 calibration remains an inherited obligation if this source is
used as a Con(ZFC) endpoint.  No result in this section fills the repository's
separate gap between an external theorem and a checked internal ZFC proof.

The executable controls test $988$ artificial Horn systems with up to three
rules on labels $0,1,2$.  They construct and verify a certificate for all
$185$ systems in that family where target~$1$ is derivable, including beta
encoding, and reject all $185$ corrupted divisor quotients.  A separate
counterfeit uses $M=1$, axiom~$3$, and rule $(1,0)\longrightarrow0$:
the accumulator $2,4,1,1$ passes the raw update but fails the divisibility
condition.  These tests support implementation checks; the unbounded theorem
is the proof above.


\section{Sources and attribution}
The archive records exact source paths and hashes in
\texttt{PROVENANCE.json}. The following references identify the mathematical
dependencies, rather than merely similar constructions.
\begin{enumerate}[leftmargin=*]
\item ProveIt, pinned wave-ten review,
\url{https://github.com/VladimirReshetnikov/ProveIt/blob/2605669a532342f153b13eb6b4a0af4b54f5fa37/Computability/DiophantineConsistency/PROPOSAL_REVIEW_WAVE10.md}.
\item Package 34, \texttt{coefficients/build\_coefficients.py},
\texttt{COEFFICIENT\_PROOF.md}, and source-34 integer/Pell calibration.
\item Package 31, the original mixed-radix coefficient tuple and its native
calibration; package 44, the distinct compact tuple used for eleven positives.
\item Package 64, \texttt{integer41/PROOF41.md}, \texttt{PROOF44.md}, and
the literal \texttt{quartic41.json}.
\item Package 68, \texttt{integer\_system/PAIR\_GUARD\_PROOF.md},
and its direct affine-product proof-certificate construction.
\item Package 69, \texttt{positive/PROOF.md}: rational Pell descent.
\item Package 70, \texttt{eleven/PROOF.md} and
\texttt{positive/section.tex}: three-root transfer and its weighted
positive transport interface.
\item Package 71, \texttt{proofs/DETERMINANT\_GUARD.md} and the finite
source-34 coefficient binding.
\item Package 65, the fixed nine-positive encoding prefix; package 59,
the inherited ratio theorem; package 71, the comparison nine-positive norm
and majorant. Their literal source files are retained in the nine-positive
component.
\item Package 75, \texttt{guards/PROOF.md} and
\texttt{guards/review/MATHEMATICAL\_REVIEW.md}: signed-power bootstrap.
\item ProveIt, \texttt{OBJECT\_ZFC.md} at the same full commit: the current
distinction between external equivalence and an internal ZFC derivation.
\end{enumerate}
Independent audits here mean separately implemented checks within the AI
collaboration; they are not external human peer review. Existing authorship,
K215, Wade, and source notices are retained. New material follows the
repository's MIT No Attribution terms, subject to inherited licences.
\end{document}
